\documentclass[11pt]{article}
\usepackage[margin=1in]{geometry}
\usepackage[T1]{fontenc}
\usepackage[utf8]{inputenc}
\usepackage{amsmath,amssymb,amsthm}
\usepackage{enumitem}

\title{ICPC World Finals 2022\\R. Zoo Management}
\author{}
\date{}

\begin{document}
\maketitle

\section*{Problem Summary}

Each move is a collection of vertex-disjoint directed cycles in the enclosure graph: every animal on such a
cycle moves one step along a tunnel, and no other animal moves. We must decide whether the initial
labeling of the vertices can be transformed into the target labeling.

\section*{Key Observations}

\begin{itemize}[leftmargin=*]
    \item A bridge can never be used in any move, because no cycle contains a bridge. So animals never
    cross bridges, and every connected component after removing bridges is independent.
    \item Inside one bridge-free component there are three qualitatively different cases:
    \begin{itemize}[leftmargin=*]
        \item the component is a single simple cycle;
        \item the component is a cactus made of several simple cycles sharing only vertices;
        \item some edge belongs to more than one cycle.
    \end{itemize}
    \item On a single simple cycle, the only possible operation is a rotation of that cycle, so the target
    sequence must be a cyclic shift of the initial sequence.
    \item If some edge belongs to more than one cycle, then one can realize arbitrary swaps and therefore
    arbitrary permutations of the vertices in that component.
    \item If the component is a cactus with more than one cycle, then all even permutations are possible.
    If at least one cycle has even length, that cycle itself is an odd permutation, so in fact all
    permutations become possible.
\end{itemize}

\section*{Algorithm}

\begin{enumerate}[leftmargin=*]
    \item Find all bridges and remove them.
    \item For each remaining connected component:
    \begin{itemize}[leftmargin=*]
        \item first compare the multisets of initial and target animal types inside the component;
        \item if the component is a simple cycle, build the cyclic order of vertices and test cyclic
        equivalence of the two label sequences;
        \item otherwise run one DFS inside the component and count, for each tree edge, how many back
        edges cover it.
    \end{itemize}
    \item If some tree edge is covered by more than one back edge, then some edge lies on multiple cycles,
    so the component allows all permutations.
    \item Otherwise the component is a cactus. Check whether one of its cycle lengths is even:
    \begin{itemize}[leftmargin=*]
        \item if yes, all permutations are possible;
        \item if not, only even permutations are possible.
    \end{itemize}
    \item In the even-only case, parity matters only when all labels in the component are distinct. Then
    compute the parity of the permutation sending the initial labeling to the target labeling.
\end{enumerate}

\section*{Correctness Proof}

We prove that the algorithm returns the correct answer.

\paragraph{Lemma 1.}
No animal can ever cross a bridge.

\paragraph{Proof.}
Each move is a union of directed cycles. A bridge belongs to no undirected cycle, so it cannot be used in
any move. Hence animals remain forever inside the same connected component of the graph after bridges
are removed. \qed

\paragraph{Lemma 2.}
If a bridge-free component is a simple cycle, then the reachable labelings are exactly the cyclic shifts of
its current labeling.

\paragraph{Proof.}
The only cycle present in the component is the whole component itself, so every move rotates the labels by
one step clockwise or counterclockwise. Repeating such moves gives exactly all cyclic shifts and nothing
else. \qed

\paragraph{Lemma 3.}
If a bridge-free component contains an edge that lies on more than one cycle, then every permutation of
its vertices is achievable.

\paragraph{Proof.}
The standard argument used in the official sketch shows that three edge-disjoint paths between two
vertices allow a transposition of arbitrary labels inside that subgraph. Combining such local swaps with the
available cycle rotations generates the full symmetric group on the component. \qed

\paragraph{Lemma 4.}
If a bridge-free component is a cactus with more than one cycle, then every even permutation is
achievable. If one of its cycles has even length, then every permutation is achievable.

\paragraph{Proof.}
For two adjacent cycles in a cactus, the commutator
\[
aba^{-1}b^{-1}
\]
acts as a 3-cycle on labels, so the alternating group is generated. Thus every even permutation is
reachable.

A cycle of length $L$ is a permutation of sign $(-1)^{L-1}$. Hence an even-length cycle is an odd
permutation. Once both all even permutations and one odd permutation are available, the full symmetric
group is generated. \qed

\paragraph{Theorem.}
The algorithm outputs \texttt{possible} exactly when the target arrangement can be achieved.

\paragraph{Proof.}
By Lemma 1, components after bridge removal are independent, so it is necessary and sufficient to solve
each one separately.

If a component is a simple cycle, Lemma 2 gives the exact criterion. Otherwise, the DFS cover counts tell
us whether some edge lies on multiple cycles. In that case Lemma 3 applies and only multiset equality is
needed. If not, the component is a cactus, so Lemma 4 tells us whether all permutations or only even
permutations are available.

In the even-only case, when some label appears at least twice, swapping two equal labels changes the
permutation parity without changing the final arrangement, so parity imposes no restriction. If all labels
are distinct, the parity of the unique permutation from the initial labeling to the target labeling is the exact
remaining condition. Therefore the algorithm is correct. \qed

\section*{Complexity Analysis}

Bridge finding and the later DFS traversals are linear. Comparing label multisets inside components
requires sorting their labels, so the total running time is $O((n+m) + n \log n)$, and the memory usage is
$O(n+m)$.

\section*{Implementation Notes}

\begin{itemize}[leftmargin=*]
    \item The code uses an iterative DFS for both bridge finding and cycle-structure analysis to avoid
    recursion-depth problems at $4 \cdot 10^5$ vertices and edges.
\end{itemize}

\end{document}
